\documentclass[13pt, a4paper]{article}

\usepackage{amsmath}
\usepackage{amssymb}
\usepackage{amsthm}
\usepackage{bussproofs}
\usepackage{fancyhdr}
\usepackage[margin=3cm]{geometry}
\usepackage{mathpartir}
\usepackage{stmaryrd}

\newtheorem{theorem}{Theorem}
\newtheorem{lemma}{Lemma}
\newtheorem{proposition}{Proposition}

\pagestyle{fancy}
\fancyhf{} % Clear all header and footer fields

% Header: left, center, right
\fancyhead[L]{Programming Languages 2}
\fancyhead[C]{Homework 05}
\fancyhead[R]{\today}
\fancyfoot[C]{\thepage}

\begin{document}

Simon Sutoris 7542170

%===============================================================================
\section*{1 Simply Typed Lambda Calculus}
%===============================================================================

\begin{prooftree}
  \AxiomC{$x: \textbf{A} \in y : \textbf{A}, x: \textbf{A}$}
  \RightLabel{\textsc{Ty-Var}}
  \UnaryInfC{$y : \textbf{A}, x: \textbf{A} \vdash x : \textbf{A}$}
  \RightLabel{\textsc{Ty-Abs}}
  \UnaryInfC{$y : \textbf{A} \vdash (\lambda x : \textbf{A}.x) : \textbf{A} \to \textbf{A} $}
  \AxiomC{$y : \textbf{A} \in y:\textbf{A}$}
  \RightLabel{\textsc{Ty-Var}}
  \UnaryInfC{$y:\textbf{A} \vdash y : \textbf{A}$}
  \RightLabel{\textsc{Ty-App}}
  \BinaryInfC{$y:\textbf{A} \vdash ((\lambda x : \textbf{A}.x)\,y) : \textbf{A}$}
\end{prooftree}


\begin{prooftree}
  \AxiomC{$y : \textbf{C} \in x : \textbf{B}, y : \textbf{C},z : \textbf{B}, w : \textbf{A}$}
  \RightLabel{\textsc{Ty-Var}}
  \UnaryInfC{$x : \textbf{B}, y : \textbf{C},z : \textbf{B}, w : \textbf{A} \vdash y : \textbf{C}$}
  \RightLabel{\textsc{Ty-Abs}}
  \UnaryInfC{$x : \textbf{B}, y : \textbf{C},z : \textbf{B} \vdash (\lambda w : \textbf{A}.y) : \textbf{A} \to \textbf{C}$}
  \RightLabel{\textsc{Ty-Abs}}
  \UnaryInfC{$x : \textbf{B}, y : \textbf{C} \vdash (\lambda z : \textbf{B}.\lambda w : \textbf{A}.y) : \textbf{B} \to (\textbf{A} \to \textbf{C})$}
  \AxiomC{$x : \textbf{B} \in x : \textbf{B}, y : \textbf{C}$}
  \RightLabel{\textsc{Ty-Var}}
  \UnaryInfC{$x : \textbf{B}, y : \textbf{C} \vdash x : \textbf{B}$}
  \RightLabel{\textsc{Ty-App}}
  \BinaryInfC{$x : \textbf{B}, y : \textbf{C} \vdash ((\lambda z : \textbf{B}.\lambda w : \textbf{A}.y)\, x) : \textbf{A} \to \textbf{C}$}
\end{prooftree}

There exist no context $\Gamma$ or type $T$ for the two latter terms $t$ in the exercise. For the first, $w$ is a free variable, and we can't deduce its type without evaluating it: Or in other words: There is no typing rule that lets us determine its type.

For the latter: $x$ must be of a type, that is both a concrete type, and can be evaluated. Both can't be simultaneously true.

\begin{prooftree}
  \AxiomC{$z : \textbf{A} \in \Gamma, x : \textbf{B}, y : \textbf{C}, z : \textbf{A} $}
  \RightLabel{\textsc{Ty-Var}}
  \UnaryInfC{$\Gamma, x : \textbf{B}, y : \textbf{C}, z : \textbf{A} \vdash z : \textbf{A}$}
  \RightLabel{\textsc{Ty-Abs}}
  \UnaryInfC{$\Gamma, x : \textbf{B}, y : \textbf{C} \vdash (\lambda z : \textbf{A}.z) : \textbf{A} \to \textbf{A}$}
  %\AxiomC{$(w\,x\,y) : \textbf{A} \in \Gamma, x : \textbf{B}, y : \textbf{C}$}
  \AxiomC{$\lightning$}
  \UnaryInfC{$\Gamma, x : \textbf{B}, y : \textbf{C} \vdash (w\,x\,y) : \textbf{A}$}
  \RightLabel{\textsc{Ty-App}}
  \BinaryInfC{$\Gamma, x : \textbf{B}, y : \textbf{C} \vdash ((\lambda z : \textbf{A}.z)\,(w\,x\,y)) : \textbf{A}$}
  \RightLabel{\textsc{Ty-Abs}}
  \UnaryInfC{$\Gamma, x : \textbf{B}\vdash \lambda y : \textbf{C}.((\lambda z : \textbf{A}.z)\,(w\,x\,y)) : \textbf{C} \to \textbf{A}$}
  \RightLabel{\textsc{Ty-Abs}}
  \UnaryInfC{$\Gamma \vdash \lambda x : \textbf{B}.\lambda y : \textbf{C}.((\lambda z : \textbf{A}.z)\,(w\,x\,y)) : \textbf{B} \to (\textbf{C} \to \textbf{A})$}
\end{prooftree}

%\begin{prooftree}
%  \AxiomC{$x : \texttt{omega} \in x : \texttt{omega}, y : \textbf{A}$}
%  \RightLabel{\textsc{Ty-Var}}
%  \UnaryInfC{$x : \texttt{omega}, y : \textbf{A} \vdash x : \texttt{omega}$}
%  \RightLabel{$\texttt{omega} = (\texttt{omega} \to \texttt{omega})$}
%  \UnaryInfC{$x : \texttt{omega}, y : \textbf{A} \vdash x : (\texttt{omega} \to \texttt{omega})$}
%  \AxiomC{$x : \texttt{omega} \in x : \texttt{omega}, y : \textbf{A}$}
%  \RightLabel{\textsc{Ty-Var}}
%  \UnaryInfC{$x : \texttt{omega}, y : \textbf{A} \vdash x : \texttt{omega}$}
%  \RightLabel{\textsc{Ty-App}}
%  \BinaryInfC{$x : \texttt{omega}, y : \textbf{A} \vdash x\,x : \texttt{omega}$}
%  \RightLabel{\textsc{Ty-Abs}}
%  \UnaryInfC{$x : \texttt{omega} \vdash \lambda y : \textbf{A}.x\,x : \textbf{A} \to \texttt{omega}$}
%\end{prooftree}

\begin{prooftree}
  \AxiomC{$\lightning$}
  \UnaryInfC{$x : \tau', y : \textbf{A} \vdash x : (\tau \to \tau')$}
  \AxiomC{$x : \tau' \in x : \tau', y : \textbf{A}$}
  \RightLabel{\textsc{Ty-Var}}
  \UnaryInfC{$x : \tau', y : \textbf{A} \vdash x : \tau'$}
  \RightLabel{\textsc{Ty-App}}
  \BinaryInfC{$x : \tau', y : \textbf{A} \vdash x\,x : \tau'$}
  \RightLabel{\textsc{Ty-Abs}}
  \UnaryInfC{$x : \tau* \vdash \lambda y : \textbf{A}.x\,x : \textbf{A} \to \tau'$}
\end{prooftree}

%===============================================================================
\section*{2 Extension with Numbers}
%===============================================================================

\begin{theorem}\label{thm:2-substitution}
Preservation of Types under Substitution (9.3.8)
\begin{prooftree}
  \AxiomC{$\Gamma, x:S\vdash t:T$}
  \AxiomC{$\Gamma \vdash s : S$}
  \BinaryInfC{$\Gamma \vdash t[s/x]:T$}
\end{prooftree}
\end{theorem}

\begin{proof}
Case \textsc{T-Var}:

If $t = x$, then $x : S \to x[s/x] = s : S \Rightarrow $ $t$ and $t[s/x]$ have the same type $S$.

If $t \ne x$, then $t \to t[s/x] : T = t : T$.

\noindent Case \textsc{T-Zero}:

$[s/x]t = \texttt{zero} : \texttt{Nat}$

\noindent Case \textsc{T-Abs}:

Assume $t = \lambda y:T_2.t_1$

and $T = T_2\to T_1$

and $\Gamma, x:S,y:T_2 \vdash t_1 :T_1$

then using permutation we get $\Gamma, y:T_2,x:S \vdash t_1 : T_1$.

Using weakening we get $\Gamma, y_T_2 \vdash s : S$

Using induction hypothesis we get $\Gamma, y:T_2 \vdash t_1[s/x]$

By \textsc{T-Abs} we get $\Gamma \vdash \lambda:T_2.t_1[s/x] : T_2\to T_1$

Thus $t[s/x] = \lambda:T_1.t_1[s/x]$.

\noindent Case \textsc{T-App}:

Assume $t = t_1\,t_2$

and $\Gamma, x:S \vdash t_y : T_2 \to T_1$

and $\Gamma, x : S \vdash t_2 : T_2$

and $T = t_1$

then by induction hypothesis $\Gamma \vdash t_1[s/x] : T_2 \to T_1$

and $\Gamma \vdash t_2[s/x] : T_2$.

Using \textsc{T-App} we get $\Gamma \vdash t_1[s/x] t_2[s/x]: T$

i.e. $\Gamma \vdash (t_1 t_2)[s/x] : T$

\noindent Case \textsc{T-Succ}:

If $t = x$, then $S = \texttt{Nat} \Rightarrow x:\texttt{Nat} \to x[s/x] = s : \texttt{Nat}$.

If $t \ne x$, then $t[s/x] = t : \texttt{Nat}$

\end{proof}

%===============================================================================

\begin{theorem}
Progress (9.3.5)\\
Let $t$ be a closed, well-typed term, that is $\vdash t:T$ for some $T$, then either $t$ is a value or there is some $t'$ with $t\to t'$.
\end{theorem}

\begin{proof}
\textsc{T-Var} does not apply, because not closed. \textsc{T-Zero}, \textsc{T-Abs}, \textsc{T-Succ} are values by definition.

For \textsc{T-App}, because of the rules axioms $\Gamma \vdash t_1 : T_{11} \to T_{12}$ and $\Gamma \vdash t_2 : T_{11}$, it follows that $t_1\,t_2 : T_{12} \to [x/t_2]t_1 : T_{12}$.
\end{proof}

%===============================================================================

\begin{theorem}
Preservation (9.3.9) \\
If $\Gamma \vdash t:T$ and $t\to t'$, then $\Gamma \vdash t' : T$.
\end{theorem}

\begin{proof}
\textsc{T-Var}, \textsc{T-Zero} and \textsc{T-Abs} don't apply, because there are no derivation rules for them.

\noindent Case \textsc{T-Succ}:

\begin{align}
& t : \texttt{Nat} \\
& t' = \texttt{succ}\,t_1 : \texttt{Nat}
\end{align}

Because $t$ and $t'$ have the same type \texttt{Nat} and $t \to t'$, \textsc{E-Succ} applies.

\noindent Case \textsc{T-App} $t_1 \to t_1'$:

By the theorem hypothesis, $t_1$ must have the same type as $t_1'$.

Using Theorem \ref{thm:2-substitution}, $t_1\,t_2$ has the same type as $t_1'\,t_2$.

\textsc{E-App1} applies.

\noindent Case \textsc{T-App} $t_1 = v_1 \land t_2\to t_2'$:

By the theorem hypothesis, $t_2$ must have the same type as $t_2'$.

Using Theorem \ref{thm:2-substitution}, $v_1\,t_2$ has the same type as $v_1\,t_2'$.

\textsc{E-App2} applies.

\noindent Case \textsc{T-App} $t_1 = (\lambda x : T.t_{12})\land t_2 = v_2$:

Due to the axioms, $t_1$ is of type $T_{11}\to T_{12}$ and $t_2 = v_2$ of type $T_{11}$.

So $v_1$ can be applied to $\lambda x : T.t_{12}$, and the reduced term is of type $T_{12}$.

\textsc{E-AppAbs} applies.

\end{proof}

%===============================================================================
\section*{3 Environments and free variables}
%===============================================================================

Exercise 1 follows from Exercise 2. Therefore I first prove Exercise 2 (Theorem \ref{thm:3-2}) and afterwards Exercise 1 (Theorem \ref{thm:3-1}).

\begin{theorem}\label{thm:3-2}
If $\Gamma \vdash t: T$, then $FV(t) \subseteq dom(\Gamma)$.
\end{theorem}

\begin{proof}

\noindent Case \textsc{T-Var}

If $t=x$, then $FV(x) = \{x\}$ and $x \in dom(\Gamma)$. Therefore $FV(t) \subseteq dom(\Gamma)$.

\noindent Case \textsc{T-Abs}

If $t=\lambda x : S. u$, then $FV(t) = FV(u) \setminus \{x\}$.

Apply theorem hypothesis to $u \Rightarrow FV(u) \subseteq dom(\Gamma) \cup \{x\}$.

Therefore: $FV(t) = FV(u) \setminus \{x\} \subseteq (dom(\Gamma) \cup \{x\}) \setminus \{x\} = dom(\Gamma)$.

\noindent Case \textsc{T-App}

Very similar to Case \textsc{T-Abs}.

If $t= t_1\,t_2$, then $FV(t) = FV(t_1) \cup FV(t_2)$.

If $t_1=\lambda x : S. u_1$, then $FV(t) = FV(u_1) \setminus \{x\}$.

Apply theorem hypothesis to $u_1 \Rightarrow FV(u_1) \subseteq dom(\Gamma) \cup \{x\}$.

Therefore: $FV(t_1) = FV(u_1) \setminus \{x\} \subseteq (dom(\Gamma) \cup \{x\}) \setminus \{x\} = dom(\Gamma)$.

Analog for $t_2=\lambda x : S. u_2$.

Therefore: $FV(t) = FV(t_1) \cup FV(t_2) \subseteq dom(\Gamma) \cup dom(\Gamma) = dom(\Gamma)$

\end{proof}


\begin{theorem}\label{thm:3-1}
If $\Gamma \vdash t: T$, and $x \notin dom(\Gamma)$, then for all values $v$ we have $t[v/x]=t$.
\end{theorem}

\begin{proof}

Position Theorem \ref{thm:3-2} : $v\in FV(t) \to v \in dom(\Gamma)$

Counterposition: $v \notin dom(\Gamma) \to v\notin FV(t)$

Therefore: $\forall v, t[v/x]=t$

\end{proof}

%===============================================================================
\section*{4 Compilation to untyped lambda calculus}
%===============================================================================

\begin{theorem}\label{thm:4-1}
$\llbracket v \rrbracket$ is a value.
\end{theorem}

\begin{proof} Case $\llbracket x \rrbracket$ and $x$ is a value

$\llbracket x \rrbracket = x$. Therefore $\llbracket x \rrbracket$ is a value.

\noindent Case $\llbracket \lambda x:T.t \rrbracket$

If $t$ is a value, then $\llbracket \lambda x:T.t \rrbracket := \lambda x.\llbracket t \rrbracket = t$ is a value.

If $t$ is not a value, then by theorem hypothesis $\llbracket t \rrbracket$ is a value.

Therefore $\lambda x.\llbracket t \rrbracket$ is in the form $\lambda x.t$ (Figure 5.3 in the Book), and thus a value.

\noindent Case $\llbracket \texttt{true} \rrbracket$

$\llbracket \texttt{true} \rrbracket := \lambda x.\lambda y.x$ is in the form $\lambda x.t$ (Figure 5.3 in the Book) and is therefore a value.

\noindent Case $\llbracket \texttt{false} \rrbracket$

Analog to case $\llbracket \texttt{true} \rrbracket$

\noindent Case $\llbracket t_1\,t_2 \rrbracket$, $\llbracket \texttt{if}\,t_1\,\texttt{then}\,t_2\,\texttt{else}\,t_3 \rrbracket$

The terms are not in the form $\llbracket v \rrbracket$ and thus the theorem does not apply.

\end{proof}

\begin{theorem}\label{thm:4-2}
If $ t\to  v$ in the calculus above, then $\llbracket t\rrbracket \to^* \llbracket v\rrbracket$ in the untyped call-by-value lambda calculus.
\end{theorem}

\begin{proof}

Case $\llbracket x \rrbracket$, $\llbracket \lambda x:T.t \rrbracket$, $\llbracket \texttt{true} \rrbracket$, $\llbracket \texttt{false} \rrbracket$

As per Theorem \ref{thm:4-1}, the terms are values, and thus Theorem \ref{thm:4-2} does not apply.

\noindent Case $\llbracket t_1\,t_2 \rrbracket$

Due to call-by-value and the theorem hypothesis, $\llbracket t_1 \rrbracket \to^* \llbracket v_1 \rrbracket$

$\llbracket v_1\,t_2 \rrbracket \to \llbracket v_1[t_2/x] \rrbracket$

Due to the \textit{progress} theorem, $v_1[t_2/x]$ either is a value, or $\to^*$ reduces it to a value.

Therefore $\llbracket v_1\,t_2 \rrbracket \to^* \llbracket v \rrbracket$

$\llbracket v_1\,t_2 \rrbracket := \llbracket v_1 \rrbracket\,\llbracket t_2 \rrbracket$

Then $\llbracket v_1 \rrbracket\,\llbracket t_2 \rrbracket \to \llbracket v_1 \rrbracket[\llbracket t_2 \rrbracket/x]$

Then the exercise allows us to assume $\llbracket v_1 \rrbracket[\llbracket t_2 \rrbracket/x] = \llbracket v_1 [t_2 /x]\rrbracket$

This reproduces the case above and we get $\llbracket v_1\,t_2 \rrbracket := \llbracket v_1 \rrbracket\,\llbracket t_2 \rrbracket \to^* \llbracket v \rrbracket$

\noindent Case $\llbracket \texttt{if}\,t_1\,\texttt{then}\,t_2\,\texttt{else}\,t_3 \rrbracket$

If $t_1 = \texttt{true}$, then \textsc{E-IfTrue} applies:

$\llbracket \texttt{if}\,t_1\,\texttt{then}\,t_2\,\texttt{else}\,t_3 \rrbracket\to \llbracket t_2 \rrbracket$

As per theorem hypothesis, $t_2 \to^* v$.

If $t_1 = \texttt{false}$, then \textsc{E-IfFalse} applies. Analog to $t_1 = \texttt{true}$.

Assuming a typo in the exercise, \textsc{E-If} is actually

\begin{prooftree}
  \AxiomC{$t_1 \to t_1'$}
  \UnaryInfC{$\texttt{if}\,t_1\,\texttt{then}\,t_2\,\texttt{else}\,t_3 \to \texttt{if}\,t_1'\,\texttt{then}\,t_2\,\texttt{else}\,t_3$}
\end{prooftree}

Then, due to the \textit{preservation} theorem, $t_1 : T \to^* t_1' : T$, $t_1$ reduces to \texttt{true} or \texttt{false}.

Then the cases above apply.

\end{proof}

%===============================================================================
\section*{5 Types and Programming Languages}
%===============================================================================

No questions.

\end{document}
